\documentclass[12pt, a4paper]{article}

\usepackage[T2A]{fontenc}
\usepackage[utf8]{inputenc}
\usepackage[russian]{babel}

\usepackage{amsmath, amssymb, amsfonts, mathtools}

\usepackage{geometry}
\geometry{left=2.5cm, right=2.5cm, top=2.5cm, bottom=2.5cm}
\usepackage{enumitem}
\usepackage{xcolor}
\usepackage{hyperref}
\hypersetup{
    colorlinks=true,
    linkcolor=blue,
    filecolor=magenta,
    urlcolor=cyan,
    pdftitle={Web},
}

% Настройка списков
\setlist[itemize]{leftmargin=1.5cm}
\setlist[enumerate]{leftmargin=1.5cm}

\begin{document}

\begin{titlepage}
    \centering
    \vspace*{5cm}
    {\Huge \textbf{Анализ социальных сетей Веб (SNA)} \par}
    \vspace{1cm}
    {\Large {Меры центральностей в Веб-графах. Примеры использования мер центральностей в прикладных задачах информационного поиска.} \par}
    \vspace{0.5cm}
    {\Large От математических моделей к прикладным архитектурам \par}
    \vspace{3cm}
    {\Large \textbf{Соснов Григорий} \par}
    \vfill
    {\large \today \par}
\end{titlepage}

\tableofcontents
\newpage

\section{Введение}

Анализ социальных сетей (Social Network Analysis, SNA) представляет собой междисциплинарную область, изучающую структуры взаимосвязей между объектами. В контексте Веб пространства объектами выступают веб-страницы, а связями — гиперссылки, формирующие ориентированный граф, известный как веб-граф. Изучение топологии веб-графа позволяет выявлять авторитетные источники, тематические сообщества и аномальные структуры, что напрямую влияет на качество информационного поиска. В условиях, когда классические модели, основанные на анализе текста, достигли насыщения, структурные меры становятся ключевым инструментом повышения релевантности и борьбы с поисковым спамом. Настоящий доклад посвящён обзору мер центральностей в веб-графах и их применению в прикладных задачах информационного поиска.


В рамках анализа веб-графов и информационного поиска можно выделить следующие ключевые задачи, решаемые с помощью мер центральностей:

\begin{enumerate}
    \item \textbf{Ранжирование веб-страниц.} Требуется упорядочить множество документов по степени их авторитетности и релевантности запросу, используя информацию о ссылочной структуре.
    \item \textbf{Поиск авторитетных источников.} Необходимо выделить страницы, которые являются «экспертами» по данной теме, даже если они не содержат точного вхождения ключевых слов.
    \item \textbf{Обнаружение тематических сообществ.} Задача состоит в разбиении веб-графа на кластеры (сообщества), внутри которых ссылки плотнее, чем между ними, что позволяет улучшить тематическую кластеризацию результатов поиска.
    \item \textbf{Выявление спама и нерелевантных страниц.} Требуется идентифицировать узлы, искусственно завышающие свою центральность через ссылочные фермы и другие манипуляции.
    \item \textbf{Персонализация и рекомендации.} На основе анализа центральностей в подграфе пользователя необходимо предлагать релевантные документы или ресурсы.
\end{enumerate}

\section{Меры центральностей в веб-графах}

Мера центральности — это функция, ставящая в соответствие каждому узлу графа числовое значение, отражающее его важность в сети. В контексте ориентированного веб-графа (где рёбра — гиперссылки) наиболее применимы следующие меры.

\subsection{Степенная центральность (Degree Centrality)}

Степенная центральность — простейшая мера, основанная на подсчёте числа связей узла. Для ориентированного графа различают входящую степень (число входящих ссылок, indegree) и исходящую степень (число исходящих ссылок, outdegree). Формально для узла $i$:

\[
C_D^{\text{in}}(i) = \sum_{j=1}^{N} a_{ji}, \quad C_D^{\text{out}}(i) = \sum_{j=1}^{N} a_{ij},
\]

где $a_{ij}$ — элемент матрицы смежности $\mathbf{A}$ графа, равный 1 при наличии ребра $i \to j$ и 0 в противном случае. Нормированная версия:

\[
C_D(i) = \frac{\deg(i)}{N-1}.
\]

В контексте Веб входящая степень традиционно рассматривается как простейший индикатор популярности страницы. Однако эта мера легко поддаётся манипуляции: создание ссылочных ферм искусственно завышает indegree. Исходящая степень характеризует «щедрость» страницы на ссылки и используется, например, для выявления страниц-каталогов. Как отмечается в обзорах, степень центральности отражает локальное лидерство мнений: высокое число прямых связей означает множество контактов в непосредственном окружении узла.

Среди значимых работ, развивающих эту меру, следует отметить исследование Бородина и соавторов «Link analysis ranking: algorithms, theory, and experiments» (2005), где проводится сравнительный анализ степенных мер и обсуждаются проблемы их устойчивости к манипуляциям. Авторы показывают, что комбинирование нескольких мер центральностей повышает качество ранжирования по сравнению с использованием только одной.

\subsection{Центральность по посредничеству (Betweenness Centrality)}

Центральность по посредничеству измеряет, насколько часто узел лежит на кратчайших путях между другими узлами. Формально:

\[
C_B(i) = \sum_{j \neq i \neq k} \frac{\sigma_{jk}(i)}{\sigma_{jk}},
\]

где $\sigma_{jk}$ — общее число кратчайших путей между узлами $j$ и $k$, а $\sigma_{jk}(i)$ — число таких путей, проходящих через узел $i$. Нормировка:

\[
C_B^{\text{norm}}(i) = \frac{2}{(N-1)(N-2)} C_B(i).
\]

В веб-графе высокая посредническая центральность указывает на страницы-«мосты», соединяющие различные тематические кластеры. Такие страницы обладают повышенной «властью» над коммуникацией: любая информация, передающаяся между кластерами через кратчайшие пути, проходит через узел с высокой посреднической центральностью. Вычисление этой меры для больших графов требует эффективных алгоритмов, таких как алгоритм Брандеса, имеющий сложность $O(NM)$ для невзвешенных графов.

Классической работой в этой области является статья Ньюмана «A measure of betweenness centrality based on random walks» (2005), где мера обобщена на случай взвешенных сетей и предложены эффективные методы обнаружения сообществ. Алгоритм Гирвана-Ньюмана, основанный на последовательном удалении рёбер с высокой посреднической центральностью, широко применяется для тематической сегментации веб-графа. В работе Чакраборти и соавторов «Mining the Web's link structure» (1999) предложены методы выявления ссылочных ферм на основе аномалий в распределении центральностей.

\subsection{Центральность по близости (Closeness Centrality)}

Центральность по близости определяется как обратная величина средней длины кратчайших путей от узла до всех остальных узлов:

\[
C_C(i) = \left( \frac{1}{N-1} \sum_{j \in V} d(i,j) \right)^{-1},
\]

где $d(i,j)$ — кратчайшее расстояние между узлами $i$ и $j$. В веб-графе эта мера показывает, насколько быстро можно достичь других страниц, что важно для оценки «навигационной» ценности документа. С точки зрения распространения влияния, узел, который легко достигает других узлов, более эффективен в распространении информации.

Однако для несвязных графов (каковым является Веб) классическое определение требует модификаций. В частности, в работе «Axioms for Centrality» (2013) Болди и Винья предлагают аксиоматический подход к определению центральностей, показывая, что closeness centrality может быть обобщена на случай несвязных графов путём добавления фиктивных рёбер или использования гармонической центральности. Практические исследования показывают, что для оценки глобального лидерства мнений closeness centrality подходит лучше, чем степенная мера, поскольку отражает не только число контактов, но и скорость достижения узлов сети.

\subsection{Центральность собственного вектора (Eigenvector Centrality)}

Центральность собственного вектора основана на идее, что узел важен, если он связан с другими важными узлами. Формально для узла $i$:

\[
C_E(i) = \frac{1}{\lambda} \sum_{j=1}^{N} a_{ij} C_E(j),
\]

где $\lambda$ — наибольшее собственное значение матрицы смежности $\mathbf{A}$. В матричной форме это уравнение записывается как $\mathbf{A}\mathbf{x} = \lambda \mathbf{x}$, где $\mathbf{x}$ — вектор центральностей.

В контексте Веб эта мера лежит в основе алгоритмов PageRank и HITS, адаптированных для веб-графов. Для scikit-learn реализован пример анализа графа ссылок внутри статей Википедии, где значения компонент первого собственного вектора используются как оценки центральности для ранжирования статей по относительной важности. В работе Кемпе, Клейнберга и Тардош «Maximizing the spread of influence through a social network» (2003) показано, что жадный алгоритм выбора seed-узлов с использованием центральности собственного вектора обеспечивает приближение к оптимальному распространению влияния.

\subsection{PageRank}

PageRank, предложенный Брином и Пейджем, моделирует случайное блуждание пользователя по ссылкам с вероятностью «телепортации» на случайную страницу. Итоговый вектор центральности удовлетворяет уравнению:

\[
\mathbf{PR} = \alpha \mathbf{A}^T \mathbf{D}^{-1} \mathbf{PR} + (1-\alpha) \frac{1}{N} \mathbf{1},
\]

где $\mathbf{A}$ — матрица смежности, $\mathbf{D}$ — диагональная матрица исходящих степеней, $\alpha$ — коэффициент затухания (обычно 0.85), а $\frac{1}{N}\mathbf{1}$ — вектор равномерного распределения телепортации. Эквивалентная итеративная форма:

\[
PR(i) = \frac{1-\alpha}{N} + \alpha \sum_{j \to i} \frac{PR(j)}{L(j)},
\]

где $L(j)$ — число исходящих ссылок страницы $j$.

PageRank является одной из наиболее влиятельных мер центральности в информационном поиске. В классической работе Брина и Пейджа «The Anatomy of a Large-Scale Hypertextual Web Search Engine» (1998) описан PageRank и архитектура поисковой системы Google, что продемонстрировало практическую ценность анализа центральностей для поиска в масштабе Веб. Современные модификации включают Personalized PageRank (для персонализации выдачи), TrustRank (для борьбы с спамом) и Topic-Sensitive PageRank (для тематического ранжирования). В работе «PageRank: Standing on the shoulders of giants» (2011) показано, что PageRank можно интерпретировать как стационарное распределение цепи Маркова, а его связь с центральностью собственного вектора позволяет строить унифицированные фреймворки для анализа ссылочных алгоритмов.

\subsection{HITS (Hyperlink-Induced Topic Search)}

Алгоритм HITS, разработанный Клейнбергом, вводит два типа центральности: \textit{authority} (авторитетность) и \textit{hub} (концентратор). Они вычисляются итеративно:

\[
    a_i = \sum_{j \to i} h_j, \quad h_i = \sum_{i \to j} a_j,
    \]
    
    где $a_i$ — authority-оценка узла $i$, $h_i$ — hub-оценка. В матричной форме:
    
    \[
        \mathbf{a} = \mathbf{L}^T \mathbf{L} \mathbf{a}, \quad \mathbf{h} = \mathbf{L} \mathbf{L}^T \mathbf{h},
        \]
        
        где $\mathbf{L}$ — матрица смежности ориентированного графа. 
        \paragraph{Определение.}Доминирующий собственный вектор — это собственный вектор матрицы или линейного оператора, который соответствует собственному значению с наибольшим модулем.
        
        Authority-вектор является доминирующим правым собственным вектором матрицы $\mathbf{L}^T\mathbf{L}$, а hub-вектор — доминирующим правым собственным вектором матрицы $\mathbf{L}\mathbf{L}^T$. Вычисление производится степенным методом с нормировкой на каждом шаге.
        
        HITS хорошо работает на подграфах, сформированных по конкретному запросу, и позволяет выделять тематические авторитеты. В работе Клейнберга «Authoritative Sources in a Hyperlinked Environment» (1999) показано, что authority-страницы содержат наиболее качественную информацию по теме, а hub-страницы являются каталогами, ссылающимися на авторитетные источники. Преимущество HITS перед PageRank состоит в том, что он предоставляет две оценки одновременно, позволяя пользователю получить как наиболее авторитетные страницы для углублённого изучения темы, так и страницы-каталоги для широкого поиска. Недостатком является более высокая уязвимость к спаму: добавление исходящих ссылок проще, чем входящих, что позволяет искусственно завышать hub-оценку.
        
        Дальнейшее развитие связано с алгоритмом SALSA, который, как и HITS, вычисляет обе оценки, но использует марковские цепи, аналогично PageRank. Это объединяет преимущества двух подходов и снижает уязвимость к манипуляциям.
        
\section{Анализ фронтиров (2025--2026)}

Ниже приведён краткий обзор новейших направлений в области мер центральностей веб-графов и их применения в информационном поиске. Материал дополняет предыдущие разделы и отражает состояние области на 2025--2026 годы.

\paragraph{Новые меры для агентного Веба.}
Классические PageRank и betweenness centrality проектировались для человеко-центричных сетей и не учитывают затухание информации при передаче между автономными агентами. В работе Samoylenko и Pribytkova (WWW Companion '26) предложена \emph{Information Spreading Betweenness Centrality} (ISBC)~--- направленная и distance-aware мера, моделирующая накопление информационных потерь вдоль кратчайших путей. Эксперименты показывают, что ISBC выявляет вершины, более эффективные по пропускной способности информации и семантической сохранности, чем PageRank и классическая betweenness centrality.

\paragraph{GNN-аппроксимация центральностей.}
Точное вычисление betweenness и closeness центральностей на больших графах вычислительно дорого. Модель \emph{CNCA-IGE} (Zou et al., KSEM 2025) использует degree centrality как входной признак и индуктивный GNN-энкодер-декодер для аппроксимации CC и BC, превосходя state-of-the-art baseline по эффективности и точности. Модель \emph{TGNN-Bet} (IEEE, 2025)~--- первый GNN-подход для аппроксимации \emph{temporal betweenness centrality} в snapshot-based временных сетях.

\paragraph{Центральности в RAG и Generative Engine Optimization (GEO).}
Переход от классического IR к Retrieval-Augmented Generation меняет структуру видимости: PageRank и Domain Authority показывают убывающую корреляцию с видимостью в LLM-экосистемах. Предложен \emph{Topological Grounding Score} (TGS)~--- фреймворк, объединяющий \emph{harmonic centrality} с семантической плотностью для количественной оценки <<grounding capacity>> домена как фактического референса для нейронного вывода. Аналогично, \emph{WebGraphMix} (Badoni et al., 2026) использует centrality scores на host-level графе Common Crawl для варьирования доли центральных и периферийных документов в pretraining mixture LLM.

\paragraph{Центральности для обнаружения спама.}
Улучшенная \emph{Entropy Centrality} (IEC) (Raj и Jaganathan, 2026) для взвешенных ориентированных графов декомпозирует граф на подграфы и анализирует прямое (1-hop) и косвенное (2--3-hop) влияние соседей, превосходя традиционные меры на датасете WEBSPAM-UK2006. Предложены также \emph{Improved H-index Centrality} и \emph{Improved Weighted H-index Centrality} для обнаружения спам-страниц.

\paragraph{Мультислойные и защищённые вычисления.}
\emph{MultiCent} (Br\"uggemann et al., 2025) обеспечивает secure и scalable вычисление центральностей на мультислойных графах, распределённых между недоверяющими сторонами, снижая асимптотическую сложность с $O(|V|^2)$ до $O(|V| \log |V|)$.

\paragraph{Спектральные фреймворки для ориентированных сетей.}
\emph{SVD Incidence Centrality} (Franco et al., arXiv 2026) предлагает унифицированный спектральный фреймворк на основе SVD матрицы инцидентности, одновременно выводя вершины и рёберные центральности через псевдоинверсию Hodge Laplacians, сохраняя направленную информацию без симметризации графа.

\paragraph{Критическая оценка.}
Большинство <<новых>> мер 2025--2026 годов~--- это либо GNN-аппроксимации уже известных центральностей (CNCA-IGE, TGNN-Bet), либо гибридные метрики, комбинирующие структурные и семантические признаки (TGS, WebGraphMix, LLM-hybrid). Фундаментально новых структурных мер, сопоставимых по влиянию с PageRank, в рассмотренных работах не предложено. Исключение~--- ISBC, но она валидирована только на синтетических данных.
        \section{Примеры использования мер центральностей в информационном поиске}
        
Приведём несколько прикладных сценариев, где меры центральностей играют ключевую роль, и рассмотрим конкретные инструменты для демонстрации.

\begin{enumerate}
\item \textbf{Ранжирование результатов поиска.} PageRank и его вариации (Personalized PageRank, TrustRank) используются для упорядочивания документов с учётом их авторитетности. Например, в системе Google PageRank долгое время был основным сигналом ранжирования.
\item \textbf{Тематический поиск и кластеризация.} Посредническая центральность позволяет выделить страницы-«мосты» между тематическими кластерами, что улучшает навигацию и тематическую сегментацию индекса. Алгоритмы обнаружения сообществ на основе центральностей применяются для построения тематических каталогов.
\item \textbf{Борьба с поисковым спамом.} Аномально высокая входящая степень в сочетании с низкой посреднической центральностью может указывать на ссылочную ферму. Системы антиспама анализируют распределение центральностей для выявления манипулятивных структур.
\item \textbf{Рекомендательные системы.} В социальных сетях и на новостных порталах центральность по собственному вектору используется для выявления влиятельных пользователей, чьи предпочтения могут быть рекомендованы другим.
\item \textbf{Персонализация поиска.} На основе подграфа, построенного по истории запросов пользователя, вычисляется персонализированный PageRank, что позволяет учитывать индивидуальные интересы.
\end{enumerate}


\newpage
\section{Заключение}

Меры центральностей в веб-графах являются мощным инструментом анализа и ранжирования, доказавшим свою эффективность в задачах информационного поиска. От классических PageRank и HITS до современных адаптивных методов — развитие этой области продолжается в направлении повышения масштабируемости, устойчивости к спаму и интеграции с семантическими моделями. Дальнейшие исследования, вероятно, будут связаны с применением графовых нейронных сетей для аппроксимации центральностей и динамическим пересчётом мер в реальном времени. Практическая демонстрация с использованием описанных инструментов и собственного анализа подграфа Википедии позволит наглядно проиллюстрировать теоретические положения доклада.

\newpage

\begin{thebibliography}{99}

\bibitem{brin1998} Brin, S., Page, L. The anatomy of a large-scale hypertextual web search engine // Computer Networks and ISDN Systems. — 1998. — Vol. 30, № 1–7. — P. 107–117.

\bibitem{kleinberg1999} Kleinberg, J. M. Authoritative sources in a hyperlinked environment // Journal of the ACM. — 1999. — Vol. 46, № 5. — P. 604–632.

\bibitem{newman2005} Newman, M. E. J. A measure of betweenness centrality based on random walks // Social Networks. — 2005. — Vol. 27, № 1. — P. 39–54.

\bibitem{kempe2003} Kempe, D., Kleinberg, J., Tardos, É. Maximizing the spread of influence through a social network // Proceedings of the 9th ACM SIGKDD International Conference on Knowledge Discovery and Data Mining. — 2003. — P. 137–146.

\bibitem{borodin2005} Borodin, A., Roberts, G. O., Rosenthal, J. S., Tsaparas, P. Link analysis ranking: algorithms, theory, and experiments // ACM Transactions on Internet Technology. — 2005. — Vol. 5, № 1. — P. 231–297.

\bibitem{chakrabarti1999} Chakrabarti, S., Dom, B., Gibson, D., Kleinberg, J., Raghavan, P., Rajagopalan, S. Mining the Web's link structure // Computer. — 1999. — Vol. 32, № 8. — P. 60–67.

\bibitem{boldi2013} Boldi, P., Vigna, S. Axioms for centrality // Internet Mathematics. — 2014. — Vol. 10, № 3–4. — P. 222–262.

\bibitem{langville2006} Langville, A. N., Meyer, C. D. Google's PageRank and Beyond: The Science of Search Engine Rankings. — Princeton University Press, 2006.

\bibitem{samoylenko2026} Samoylenko, I., Pribytkova, E. How to Measure Node Importance in the Agentic Web // Companion Proceedings of the ACM Web Conference 2026 (WWW Companion '26). — 2026. — DOI: 10.1145/3774905.3795094.

\bibitem{zou2025} Zou, Y., Li, T., et al. Node Centrality Approximation in Complex Networks via Inductive Graph Neural Networks // Knowledge Science, Engineering and Management (KSEM 2025). — Springer, 2025. — DOI: 10.1007/978-981-95-3055-7\_3.

\bibitem{sadhu2025} Sadhu, S., Kumari, K., et al. TGNN-Bet: Approximation of Temporal Betweenness Centrality using Temporal Graph Neural Network // 2025 17th International Conference on COMmunication Systems and NETworks (COMSNETS). — IEEE, 2025. — P. 911–915. — DOI: 10.1109/COMSNETS63942.2025.10885618.

\bibitem{mosquete2026} Sánchez Mosquete, J. Topological SEO and GEO: Empirical Validation of Structural Efficiency and the Topological Grounding Score for Generative Engine Optimization. — Zenodo, 2026. — DOI: 10.5281/zenodo.18552471.

\bibitem{badoni2026} Badoni, V., Chen, D., Wang, X. Hubs or Fringes: Pretraining Data Selection via Web Graph Centrality // ArXiv. — 2026. — Vol. abs/2606.11499.

\bibitem{raj2026} Raj, D. S., Jaganathan, B. An improved entropy centrality for detection of spam webpages // Cogent Engineering. — 2026. — Vol. 13, № 1. — DOI: 10.1080/23311916.2026.2627086.

\bibitem{raj2026hindex} Raj, D. S., Jaganathan, B. Detection of spam webpages using a new improved H-index centrality and weighted H-index centrality // New Review of Hypermedia and Multimedia. — 2026.

\bibitem{bruggemann2025} Brüggemann, A., Koti, N., Kukkala, V. B., Schneider, T. MultiCent: Secure and Scalable Computation of Centrality Measures on Multilayer Graphs // Proceedings on Privacy Enhancing Technologies Symposium (PoPETs). — 2025. — Cryptology ePrint Archive, Paper 2025/652.

\bibitem{franco2026} Franco, J. L., Peron, T., Dal Col, A., Petronetto, F., Verri, F. A. N., Tokuda, E. K., Nonato, L. G. SVD Incidence Centrality: A Unified Spectral Framework for Node and Edge Analysis in Directed Networks and Hypergraphs // ArXiv. — 2026. — Vol. abs/2602.15736.

\end{thebibliography}

\end{document}